\documentclass[UTF8,11pt,fontset=none]{ctexart}
\usepackage[a4paper,margin=24mm]{geometry}
\usepackage{amsmath,amssymb,mathtools}
\usepackage{enumitem}
\setCJKmainfont{Songti SC}
\setmainfont{Times New Roman}
\setlength{\parindent}{0pt}
\setlist[enumerate]{leftmargin=2.3em,itemsep=.85em,topsep=.5em}
\pagestyle{empty}
\begin{document}
\begin{center}
  {\Large\bfseries 浙江大学数学科学学院 2023 年直博生招生考试试题}
\end{center}
\vspace{1em}

\section*{数学分析}
\begin{enumerate}
\item 计算 $\displaystyle\sum_{n=0}^{\infty}\frac{1}{(2n)!}$。
\item 设 $f(x)=\displaystyle\sum_{k=1}^{n}a_k\sin kx$ 满足 $|f(x)|\le |\sin x|$。求证：
\[
\left|\sum_{k=1}^{n}ka_k\right|\le 1.
\]
\item 设 $y=y(x)$ 由方程 $y+2y^2+y^3=e^{-x}+x-1$ 确定，且
\[
y=Ax^2+Bx^3+o(x^3).
\]
求 $A,B$ 的值。
\item 设 $f(x)$ 在 $\mathbb R$ 上满足 Lipschitz 条件，证明 $f(x)$ 在 $\mathbb R$ 上几乎处处可导。
\item 已知 $f\in C[0,1]$，$A=\displaystyle\int_0^1f(x)\,dx$，
\[
h_n(x)=
\begin{cases}
nx,&0\le x\le\dfrac1{3n},\\[.35em]
\dfrac12-\dfrac12nx,&\dfrac1{3n}\le x\le\dfrac1n,\\[.35em]
0,&\text{其他},
\end{cases}
\qquad
S_n(x)=\sum_{k=0}^{n-1}h_n\!\left(x-\frac{k}{n}\right).
\]
求 $\displaystyle\lim_{n\to\infty}\int_0^1 f(x)S_n(x)\,dx$（用 $A$ 表示）。
\end{enumerate}

\clearpage
\section*{高等代数}
\begin{enumerate}
\item 已知矩阵
\[
A=\begin{pmatrix}
0&0&0&0&-7\\
-1&0&0&0&5\\
0&-1&0&0&0\\
0&0&-1&0&0\\
0&0&0&-1&0
\end{pmatrix}.
\]
求 $A$ 的极小多项式，并证明在复数域上该矩阵可以对角化。
\item 设矩阵 $A=P^{\mathrm T}P$，$B=-Q^{\mathrm T}Q$，其中 $P,Q$ 都是数域上的满秩方阵。若方阵 $C$ 满足 $AC=CB$，求矩阵 $C$。
\item 假设
\[
\alpha=(a_1,\ldots,a_n)^{\mathrm T},\quad
\beta=(b_1,\ldots,b_n)^{\mathrm T},\quad
\gamma=(c_1,\ldots,c_n)^{\mathrm T},\quad
\delta=(d_1,\ldots,d_n)^{\mathrm T}
\]
都是 $\mathbb R^n$ 中的列向量，计算行列式
\[
\left|2I_n+\alpha\beta^{\mathrm T}+\gamma\delta^{\mathrm T}\right|.
\]
\end{enumerate}
\end{document}
